%=============================================================
%  Module I.4 / L07 — Stark Effect
%  Series I > Module I.4 > Lecture L07
%  Duration: 90 minutes  |  All tiers (T1–T5)
%=============================================================
\renewcommand{\lectureCode}{I.4 / L07}

\section*{L07 — Stark Effect}
\label{sec:L07}
\addcontentsline{toc}{section}{L07 — Stark Effect}

%--- Lecture header
\begin{tcolorbox}[lecturebox, title={Module I.4 / L07 — Metadata}]
\begin{tabular}{@{}ll@{}}
\textbf{Series:}    & I — Introductory Quantum Mechanics \\
\textbf{Module:}    & I.4 — Paradigm Physical Systems \\
\textbf{Lecture:}   & L07 — Stark Effect \\
\textbf{Duration:}  & 90 minutes \\
\textbf{Tiers:}     & T1 (HS) · T2 (BegUG) · T3 (AdvUG) · T4 (MSc) · T5 (PhD) \\
\textbf{Preceding:} & I.4 / L06 — Zeeman Effect \\
\textbf{Following:} & I.4 / L08 — Spin–Orbit Coupling and Fine Structure \\
\textbf{Prerequisites:} & Modules I.1–I.3 complete \\
\end{tabular}

\medskip
\textbf{Key concepts:} Parity selection rule, linear Stark H n=2, quadratic Stark polarisability, parabolic coordinates
\end{tcolorbox}

%------------------------------------------------------------
\subsection*{L07.0 — Learning Objectives (per tier)}
\label{sec:L07.0}

\tiertag{tier1col}{Tier 1 — HS}\quad
Identify the main physical phenomenon studied in this lecture; describe it in
non-mathematical terms; state one real-world application.

\tiertag{tier2col}{Tier 2 — BegUG}\quad
Apply the key equations to compute numerical results; verify consistency with
known limits; translate between physical and mathematical descriptions.

\tiertag{tier3col}{Tier 3 — AdvUG}\quad
Derive central results from first principles; prove key identities; identify
the mathematical structure underlying the physical phenomenon.

\tiertag{tier4col}{Tier 4 — MSc}\quad
Analyse formal extensions; connect to symmetry principles; generalise to
related systems; identify the role of this lecture in the broader programme.

\tiertag{tier5col}{Tier 5 — PhD}\quad
Establish rigorous mathematical foundations; identify open research problems;
connect to modern literature; critique standard textbook treatments.

%--- Lecture content -----------------------------------------

\subsection*{L07.1 — Why First-Order Stark Vanishes}
For non-degenerate parity eigenstates: $\langle nlm|\hat{z}|nlm\rangle = 0$
by parity ($\hat{z}$ is odd under $\mathbf{r}\to-\mathbf{r}$).
Exception: degenerate $n=2$ states of hydrogen.

\subsection*{L07.2 — Linear Stark Effect in H ($n=2$)}
Four degenerate states $\{2s,2p_{-1},2p_0,2p_1\}$.
Non-zero matrix element: $\langle 2s|\hat{H}'|2p_0\rangle = -3eEa_0$.
Eigenvalues of the $4\times4$ matrix: $\{-3eEa_0,0,0,+3eEa_0\}$.

\subsection*{L07.3 — Quadratic Stark and Polarisability}
\begin{tcolorbox}[lecturebox, title={Polarisability}]
\begin{equation}
  E^{(2)}_0 = \sum_{n\neq0}\frac{|\langle n|\hat{H}'|0\rangle|^2}{E_0-E_n}
  = -\frac{9}{4}a_0^3 E^2,
  \quad \alpha = \frac{9}{2}a_0^3.
  \label{eq:L07.1}
\end{equation}
\end{tcolorbox}

\subsection*{L07.4 — Parabolic Coordinates}
Exact separation in $(\xi,\eta,\phi)$; quantum numbers $n_1,n_2,m$;
first-order energy: $E^{(1)}_{n_1n_2m} = -\frac{3}{2}neEa_0(n_1-n_2)$.


%------------------------------------------------------------
\subsection*{L07.C — Connections}
\label{sec:L07.C}
\textbf{Backward:} I.4/L06 — Zeeman Effect and Modules I.1–I.3.
\textbf{Forward:} I.4/L08 — Spin–Orbit Coupling and Fine Structure builds directly on the formalism developed here.
\textbf{Cross-module:} Series~II modules extend these results to many-body and
advanced settings; Series~III applies them to molecular electronic structure.

%=============================================================
%  ASSESSMENT BUNDLE — L07
%=============================================================
\newpage
\section*{L07 — Assessment Artifact Bundle}

\begin{tcolorbox}[ugconceptbox, title={SCQU · L07 · Tiers 1–3 · 15–20 min}]
\textit{Ten simple concept questions (mix MC/TF/short) targeting the core results
of Stark Effect. See artifact file \texttt{L07_artifacts.tex} for full questions.}

\textbf{Rubric:} 2 pts each (1.5 correct + 0.5 justification). Total: 20 pts.
\end{tcolorbox}

\begin{tcolorbox}[ugconceptbox, title={CCQU · L07 · Tiers 2–3 · 25–35 min}]
\textit{Ten complex concept questions: ``explain why,'' ``what would change if,''
``identify the flaw.'' Full questions in \texttt{L07_artifacts.tex}.}

\textbf{Rubric:} 2 pts each (1 key point + 1 argument). Total: 20 pts.
\end{tcolorbox}

\begin{tcolorbox}[ugprobbox, title={SPSU · L07 · Tiers 1–3 · 60–90 min}]
\textit{Ten problems (Part A: mechanics, Part B: interpretation, Part C: translation).
Full problems in \texttt{L07_artifacts.tex}.}

\textbf{Rubric:} Result 60\%, Method 30\%, Interpretation 10\%.
\end{tcolorbox}

\begin{tcolorbox}[ugprobbox, title={CPSU · L07 · Tier 3 · 3–6 hrs (4-layer)}]
\textit{Five multi-part derivation problems with four-layer scaffolding.
Layers 1–4 in \texttt{L07_ex01\_problem.tex}, \texttt{hints}, \texttt{adv\_hints}, \texttt{solution}.}

\textbf{Rubric:} Assumptions 15\%, Proof structure 45\%, Algebra 25\%, Insight 15\%.
\end{tcolorbox}

\begin{tcolorbox}[ugprobbox, title={SPJU + CPJU · L07 · Tiers 2–3}]
\textit{Simple project (2–6 hrs): structured computational/analytical task.
Complex project (8–15 hrs): open-ended synthesis.
Full specs in \texttt{L07_artifacts.tex}.}
\end{tcolorbox}

\begin{tcolorbox}[gradconceptbox, title={SCQG + CCQG · L07 · Tiers 4–5}]
\textit{Ten simple + ten complex concept questions at graduate level (rigour, domain
conditions, named theorems). Full questions in \texttt{L07_artifacts.tex}.}
\end{tcolorbox}

\begin{tcolorbox}[gradprobbox, title={SPSG + CPSG · L07 · Tiers 4–5}]
\textit{Ten simple + five complex graduate problems.
CPSG with four-layer format. See \texttt{L07_artifacts.tex}.}
\end{tcolorbox}

\begin{tcolorbox}[gradprobbox, title={SPJG (×5) + CPJG (×5) · L07 · Tiers 4–5}]
\textit{Five simple (3–8 hrs each) and five complex (10–20 hrs each) graduate projects.
Full specs in \texttt{L07_artifacts.tex}.}
\end{tcolorbox}

\begin{tcolorbox}[researchbox, title={WRQ (×5) + ORQ (×5) · L07 · Tiers 4–5}]
\textit{Five well-defined research questions (5–10 hrs each) and five open-ended
research questions (10–20 hrs each). Full prompts in \texttt{L07_artifacts.tex}.}
\end{tcolorbox}
